\documentclass[10pt,a4paper]{amsart}
\usepackage{amsmath,amssymb,amsthm,graphicx}
\usepackage{hyperref}
\usepackage{tikz}
%\usepackage{xcolor}
\usepackage{pict2e}
\let\arc\relax
%\usepackage{harpoon}
\usepackage{curves}
%\usepackage{chngcntr}
%\usepackage{apptools}
\usepackage{enumerate}

%\usepackage{verbatim}
%\usepackage[nospace,noadjust]{cite}
\usepackage{accents}
\usetikzlibrary{calc,decorations.markings}
\unitlength1cm

%%%%%%%%%%%%%%FONTS%%%%%%%%%%%%%%%%%%%%%%%%%%%%%%
\newcommand{\R}{\mathbb{R}}
\newcommand{\Z}{\mathbb{Z}}
\newcommand{\N}{\mathbb{N}}
\newcommand{\C}{\mathbb{C}}
\newcommand{\Id}{I}

%%%%%%%%%%%%%%MATHCAL%%%%%%%%%%%%%%%%%%%%%%%
\newcommand{\cK}{\mathcal{K}}
\newcommand{\cE}{\mathcal{E}}
\newcommand{\cR}{\mathcal{R}}
\newcommand{\cW}{\mathcal{W}}

%%%%%%%%%%%%%%%%%%ABBRS%%%%%%%%%%%%%%%%%%%%%%%%%%%%%

%%%%%%% OBJECTS SPECIFIC TO THIS PAPER%%%%%%%%%

\newcommand{\hw}{\widehat w}
\newcommand{\tw}{\widetilde w}

\newcommand{\bQ}{ Q}
\newcommand{\hQ}{{\widehat Q}}
\newcommand{\tQ}{{\widetilde Q}}

\newcommand{\sbQ}{Q^\star}
\newcommand{\shQ}{{\widehat Q^\star}}
\newcommand{\stQ}{{\widetilde Q^\star}}

\newcommand{\smu}{\mu^\star}
\newcommand{\shmu}{\widehat \mu^\star}
\newcommand{\stmu}{\widetilde \mu^\star}

\newcommand{\bmu}{{\mu}}
\newcommand{\hmu}{{\widehat \mu}}
\newcommand{\tmu}{{\widetilde \mu}}


\newcommand{\bP}{P}
\newcommand{\hP}{{\widehat P}}
\newcommand{\tP}{{\widetilde P}}

\newcommand{\sbP}{P^\star}
\newcommand{\shP}{{\widehat P^\star}}
\newcommand{\stP}{{\widetilde P^\star}}

\newcommand{\hR}{{\widehat R}}

\newcommand{\bv}{v}
\newcommand{\hv}{{\widehat v}}
\newcommand{\tv}{{\widetilde v}}

\newcommand{\sbv}{v^\star}
\newcommand{\shv}{{\widehat v^\star}}
\newcommand{\stv}{{\widetilde v^\star}}

\newcommand{\stu}{{\widetilde u^\star}}

\newcommand{\hF}{\widehat F}
\newcommand{\hW}{\widehat W}

\newcommand{\Fw}{\frac{F^2}{w}}
\newcommand{\hFw}{\frac{\hF^2}{\hw}}

\newcommand{\Cau}{\mathcal C_{\Sigma}}
\newcommand{\Cauv}{\mathcal C_{v}}
\newcommand{\Cauhv}{\mathcal C_{\hv}}
\newcommand{\Cautv}{\mathcal C_{\tv}}

\newcommand{\sCau}{\mathcal C_{\Sigma^\star}}
\newcommand{\sCauv}{\mathcal C_{\bv^\star}}
\newcommand{\sCauhv}{\mathcal C_{\hv^\star}}
\newcommand{\sCautv}{\mathcal C_{\stv}}

\newcommand{\Ltwo}{L^2(\Sigma)}
\newcommand{\sLtwo}{L^2(\Sigma^\star)}

\newcommand{\Sigl}{\Sigma^\ell}
\newcommand{\Sigr}{\Sigma^r}

%%%%%%%%%%%%%%%%OTHER SHORTCUTS%%%%%%%%%%%
\newcommand{\nn}{\nonumber}
\newcommand{\ul}{\underline}
\newcommand{\ol}{\overline}
\newcommand{\pa}{\partial}
\newcommand{\ti}{\tilde}
\newcommand{\wti}{\widetilde}
\newcommand{\wha}{\widehat}
\newcommand{\norm}[1]{\lVert#1 \rVert}
\newcommand{\abs}[1]{\lvert#1 \rvert}
%\newcommand{\Id}{\mathbb{I}}
\newcommand{\I}{\mathrm{i}}
\newcommand{\E}{\mathrm{e}}
\newcommand{\re}{\mathop{\mathrm{Re}}}
\newcommand{\im}{\mathop{\mathrm{Im}}}
\newcommand{\eps}{\varepsilon}
\newcommand{\dist}{\mathop{\mathrm{dist}}}

\newtheorem{Theorem}{Theorem}[section]
\newtheorem{Corollary}[Theorem]{Corollary}
\newtheorem{Lemma}[Theorem]{Lemma}
\newtheorem{Proposition}[Theorem]{Proposition}
 { \theoremstyle{definition}
\newtheorem{Definition}[Theorem]{Definition}
\newtheorem{Note}[Theorem]{Note}
\newtheorem{Example}[Theorem]{Example}
\newtheorem{Remark}[Theorem]{Remark} }

\numberwithin{equation}{section}
\usepackage{cleveref}
\crefname{originalequation}{original equation}{original equations}
\Crefname{originalequation}{Original equation}{Original equations}
\makeatletter
\newlabel{Tjump-1}{{internal}{source}{}{sourcepointer-1}{}}
\newlabel{Tjump-1@cref}{{[originalequation][0][]internal}{[1][0][]source}{}{}{}}
\makeatother
\crefname{originalreference}{original reference}{original references}
\Crefname{originalreference}{Original reference}{Original references}
\makeatletter
\newlabel{lensJumpContour}{{internal}{source}{}{sourcepointer-2}{}}
\newlabel{lensJumpContour@cref}{{[originalreference][0][]internal}{[1][0][]source}{}{}{}}
\makeatother
\crefname{originalequation}{original equation}{original equations}
\Crefname{originalequation}{Original equation}{Original equations}
\makeatletter
\newlabel{TQRel}{{internal}{source}{}{sourcepointer-3}{}}
\newlabel{TQRel@cref}{{[originalequation][0][]internal}{[1][0][]source}{}{}{}}
\makeatother
\crefname{originalequation}{original equation}{original equations}
\Crefname{originalequation}{Original equation}{Original equations}
\makeatletter
\newlabel{3/2decay}{{internal}{source}{}{sourcepointer-4}{}}
\newlabel{3/2decay@cref}{{[originalequation][0][]internal}{[1][0][]source}{}{}{}}
\makeatother
\title{The difference of the first recurrence matrices as a Cauchy integral}
\author{Deift, Percy, Piorkowski, Mateusz}
\date{doi:10.3842/SIGMA.2024.004; CC BY 4.0}
\begin{document}
\maketitle
\setcounter{section}{3}
\setcounter{Theorem}{4}
\noindent\textit{Editorial note.} All delivered mathematical text is the author's own, copied from the pinned SIGMA article source; no statement or proof step was rewritten, reordered or removed. The counted statement is the proposition environment at source lines 916-921 and the counted proof is the proof environment at source lines 922-992. The environment itself carries no label; the anchor is the label of the first displayed alignment inside the statement, exactly as the intake census pairs it.\par\medskip
\noindent\textit{Editorial note.} The intake census pairs this statement with a proof environment that the source places after two further statements of the same block; reading the source shows that the proof derives exactly the formula of the counted statement and that the two statements placed before it are quoted results of the cited literature, so the counted statement has its own proof.\par\medskip
\noindent\textit{Editorial note.} The proof cites published references of the author's bibliography by key; the cited entries are copied line for line from the author's own bibliography into a retained bibliography fragment, and hyperref is added because the publisher class, not the author preamble, supplied the link command.\par\medskip
\noindent\textit{Editorial note.} The retained preamble of this unit was generated mechanically for this delivery: it copies the author's own package, numbering, macro and theorem declaration block from the pinned source (source lines 9-133), with the publisher class and page setup lines removed.\par\medskip
\noindent\textit{Allowed standard prerequisites (editorial).}\begin{itemize}
\item The Riemann-Hilbert problem for the two weights and its jump matrices. same article, the displays at source lines 674 and 745 and the relation at source line 761. Fixes the jump matrices of the two solutions and the relations between them, which are what the auxiliary matrix of the proof is built from.
\item The Fokas-Its-Kitaev problem for the companion weight and its solution. same article, the display at source line 745. Provides the model solution with respect to which the behaviour of the auxiliary matrix is controlled.
\item The Sokhotski-Plemelj formula for a matrix valued function with a jump on a contour. standard Riemann-Hilbert material used at proof lines 983-989. Recovers the auxiliary matrix from its jump and its normalization at infinity, which is the step that produces the integral formula.
\end{itemize}\medskip
\ifdefined\hypertarget\hypertarget{sourcepointer-1}{}\fi
\noindent\textit{Source reference.} Original equation internal: sources/sigma-2024-004-recurrence-coefficients-for-orthogonal-polyn/pinned-original-source.tex, source lines 674--678. The counted proof refers to that the author's displayed jump matrix of the first solution; it is the author's own earlier display, definition, figure or statement in the same article and the whole argument of the counted span is retained verbatim.\par\medskip
\ifdefined\hypertarget\hypertarget{sourcepointer-2}{}\fi
\noindent\textit{Source reference.} Original reference internal: sources/sigma-2024-004-recurrence-coefficients-for-orthogonal-polyn/pinned-original-source.tex, source lines 745--749. The counted proof refers to that the author's displayed model solution of the companion Riemann-Hilbert problem; it is the author's own earlier display, definition, figure or statement in the same article and the whole argument of the counted span is retained verbatim.\par\medskip
\ifdefined\hypertarget\hypertarget{sourcepointer-3}{}\fi
\noindent\textit{Source reference.} Original equation internal: sources/sigma-2024-004-recurrence-coefficients-for-orthogonal-polyn/pinned-original-source.tex, source lines 761--765. The counted proof refers to that the author's displayed relation between the two solutions; it is the author's own earlier display, definition, figure or statement in the same article and the whole argument of the counted span is retained verbatim.\par\medskip
\ifdefined\hypertarget\hypertarget{sourcepointer-4}{}\fi
\noindent\textit{Source reference.} Original equation internal: sources/sigma-2024-004-recurrence-coefficients-for-orthogonal-polyn/pinned-original-source.tex, source lines 541--545. The counted proof refers to that the author's displayed decay estimate for the exponential factor; it is the author's own earlier display, definition, figure or statement in the same article and the whole argument of the counted span is retained verbatim.\par\medskip
\begin{Proposition}[{\cite[Proposition~4.9]{CD}}]
The following formula holds:
\begin{align}\label{QIntForm}
 \bQ^{(n)}_1- \hQ^{(n)}_1 = -\frac{1}{2\pi\I}\int_\Sigma \bQ^{(n)}_-(s)\big(\bv^{(n)}(s)-\hv^{(n)}(s)\big)\big[\hQ^{(n)}_-(s)\big]^{-1} {\rm d}s.
\end{align}
\end{Proposition}

\begin{proof}
In the following, we drop the subscript $(n)$ for better readability, i.e., write $\bQ_1$ for $\bQ^{(n)}_1$ and so on. Let us define the matrix-valued function $X(z) = \bQ(z) \big[\hQ(z)\big]^{-1}$ for $z \in \C \setminus \Sigma$. Note that
\begin{align}\label{Rformula}
 X(z) = \Id + \frac{\bQ_1-\hQ_1}{z} + O\left(\frac{1}{z^2}\right) \qquad \text{as } z \to \infty.
\end{align}
Moreover, $X$ satisfies the jump condition
\begin{align*}
 X_+(s) = X_-(s)v_X(s), \qquad s\in \Sigma,
\end{align*}
where $v_X = \hQ_- \bv \hv^{-1} \hQ_-^{-1}$. We claim that $X_\pm \in L^1(\Sigma)$. To see this let us first introduce the analog of $T$ for the weight function $\hw$:
\begin{align*}
 \widehat T(z) = \left(\frac{2}{\hF_\infty}\right)^{n\sigma_3} \widehat Y(z) \left(\frac{\hF(z)}{\phi^n(z)}\right)^{\sigma_3}.
\end{align*}
Here $\hat Y$ is the solution to the Fokas--Its--Kitaev problem for the weight function $\hw$. Then for~${z \in \Omega_0}$ (cf.\ Figure~\ref{lensJumpContour}), we have
\begin{align*}%\label{RTT}
 X(z) = T(z)\big[\widehat T(z)\big]^{-1} = O(1), \qquad z\in \Omega_0, \quad z\to-1,
\end{align*}
where we have used \eqref{Tjump-1} and its analog (see \cite[Section~2]{KMVV})
\begin{align}\label{hTjump-1}
 \widehat T(z) = \begin{pmatrix}
 O\big(|z+1|^{1/2}\big) & O\big(|z+1|^{-1/2}\big)
 \\
 O\big(|z+1|^{1/2}\big) & O\big(|z+1|^{-1/2}\big)
 \end{pmatrix}, \qquad z \to -1.
\end{align}
Note that the analog of \eqref{TQRel} remains valid
\begin{align*}
 \hQ(z) =
\begin{cases}
\widehat T(z), & z \in \Omega_0,
\\
\widehat T(z)
\begin{pmatrix}
1 & 0
\\
-\frac{\hF^2}{\hw}(z) \phi^{-2n}(z) & 1
\end{pmatrix},
& z \in \Omega_1,\vspace{1mm}
\\
\widehat T(z)
\begin{pmatrix}
1 & 0
\\
\frac{\hF^2}{\hw}(z) \phi^{-2n}(z) & 1
\end{pmatrix},& z \in \Omega_2.
\end{cases}
\end{align*}
Thus for $z\in \Omega_1\cup \Omega_2$, we have
\begin{align*}
 &X(z) = T(z)
 \begin{pmatrix}
 1 & 0
 \\
 \mp\big(\frac{F^2}{w}(z)-\frac{\hF^2}{\hw}(z)\big) \phi^{-2n}(z) & 1
 \end{pmatrix}
 \big[\widehat T(z)\big]^{-1} = O(1),
 \\
 &z\in \Omega_1\cup\Omega_2, \qquad z\to-1,
\end{align*}
where the $+$, resp.\ $-$ sign refer to $\Omega_1$ and $\Omega_2$. Apart from \eqref{Tjump-1} and \eqref{hTjump-1}, we have used in addition \eqref{3/2decay}. As $X$ can have only logarithmic singularities for $z \to 1$ and otherwise the limits to the contour are analytic, it follows that $X_\pm \in L^1(\Sigma)$.

Next, we can use the Sokhotski--Plemelj theorem to conclude that $X$ can be represented as
\begin{align*}
 X(z) &= \Id + \frac{1}{2\pi \I}\int_\Sigma \frac{X_+(s)-X_-(s)}{s-z} {\rm d}s
 \\
 &= \Id + \frac{1}{2\pi \I}\int_\Sigma \frac{Q_-(s)\big(\bv(s)[\hv(s)]^{-1}-\Id\big)\big[\hQ_-(s)\big]^{-1}}{s-z} {\rm d}s
 \\
 &= \Id + \frac{1}{2\pi \I}\int_\Sigma \frac{Q_-(s)(\bv(s)-\hv(s))[\hQ_-(s)]^{-1}}{s-z} {\rm d}s,
\end{align*}
where we made use of the particular form of the jump matrices $\bv$ and $\hv$. This representation together with \eqref{Rformula} implies \eqref{QIntForm}.
\end{proof}

\begin{thebibliography}{99}
\bibitem{CD}
Conway T.O., Deift P., Asymptotics of polynomials orthogonal with respect to a
 logarithmic weight, \href{https://doi.org/10.3842/SIGMA.2018.056}{\textit{SIGMA}} \textbf{14} (2018), 056, 66~pages,
 \href{https://arxiv.org/abs/1711.01590}{arXiv:1711.01590}.


\bibitem{KMVV}
Kuijlaars A.B.J., McLaughlin K.T.-R., Van~Assche W., Vanlessen M., The
 {R}iemann--{H}ilbert approach to strong asymptotics for orthogonal
 polynomials on {$[-1,1]$}, \href{https://doi.org/10.1016/j.aim.2003.08.015}{\textit{Adv. Math.}} \textbf{188} (2004), 337--398,
 \href{https://arxiv.org/abs/math.CA/0111252}{arXiv:math.CA/0111252}.
\end{thebibliography}
\end{document}
